\documentclass[aspectratio=169,10pt]{beamer}

\usetheme{Madrid}
\usecolortheme{default}
\setbeamertemplate{navigation symbols}{}
\setbeamertemplate{footline}[frame number]

\usepackage{amsmath,amssymb,amsthm,booktabs,mathtools}
\usepackage{graphicx}
\usepackage{hyperref}
\usepackage{tikz}

% Semantic colors from the beamer skill
\definecolor{positive}{HTML}{0173B2}       % blue
\definecolor{negative}{HTML}{DE8F05}       % orange
\definecolor{emphasis}{HTML}{029E73}       % green

\newcommand{\pos}[1]{\textcolor{positive}{\textbf{#1}}}
\newcommand{\con}[1]{\textcolor{negative}{\textbf{#1}}}
\newcommand{\HL}[1]{\textcolor{emphasis}{\textbf{#1}}}

\title[CodeGraph: Structural Retrieval]{CodeGraph: Repository-Level Structural Retrieval}
\subtitle{Deterministic repository retrieval for coding agents under fixed context budgets}
\author[Toma Alexandru-Mihai]{Toma Alexandru-Mihai}
\institute[Babes-Bolyai University]{
  Computer Science Department\\
  Babes-Bolyai University
}
\date[SCSS 2026]{SCSS Conference 2026}

\begin{document}

% Slide 1
\begin{frame}
  \titlepage
\end{frame}

% Slide 2
\begin{frame}{Coding agents fail when issue text and identifiers do not line up}
  \begin{columns}[T]
    \column{0.48\textwidth}
      \begin{block}{The Context Requirement}
        Autonomous coding agents need the \pos{exact, relevant code files} placed into their limited context window to produce valid patches.
      \end{block}
      \vspace{10pt}
      \textbf{The High Stakes:}\\
      Without the correct context, even the most advanced language models will \con{hallucinate or fail} to resolve the issue.

    \column{0.48\textwidth}
      \begin{alertblock}{The Vocabulary Gap}
        Issue descriptions are written in \textit{natural language}, which often lacks the specific \textit{code identifiers} present in the repository.
      \end{alertblock}
  \end{columns}
\end{frame}

% Slide 3
\begin{frame}{Sparse and dense retrieval miss structural reachability}
  \begin{columns}[T]
    \column{0.32\textwidth}
      \begin{alertblock}{Sparse (e.g., BM25)}
        \begin{itemize}
          \item Requires exact lexical overlap.
          \item \con{Misses dependencies} that share no vocabulary.
        \end{itemize}
      \end{alertblock}
      
    \column{0.32\textwidth}
      \begin{alertblock}{Dense (e.g., CodeBERT)}
        \begin{itemize}
          \item Captures semantic similarity.
          \item \con{Ignores explicit execution edges} in the codebase.
        \end{itemize}
      \end{alertblock}

    \column{0.32\textwidth}
      \begin{exampleblock}{CodeGraph (Ours)}
        \begin{itemize}
          \item \pos{Navigates reachability}.
          \item Ranks through actual imports, calls, and inheritance.
        \end{itemize}
      \end{exampleblock}
  \end{columns}
  \vspace{15pt}
  \centering
  \textbf{Code is a graph, not a bag of tokens.} Retrieving it as flat text ignores how the program actually runs and interacts.
\end{frame}

% Slide 4
\begin{frame}{The query chooses where relevance starts; the graph decides the flow}
  \begin{columns}[T]
    \column{0.48\textwidth}
      \begin{block}{1. Lexical Seeding}
        Natural-language issue text is used strictly to identify plausible \pos{entry points}. This places initial probability mass onto specific nodes in the graph, independent of structural links.
      \end{block}
      
    \column{0.48\textwidth}
      \begin{exampleblock}{2. Structural Propagation}
        The final relevance ranking is computed \pos{entirely by graph topology}. Relevance expands outward from the seeds along true repository dependencies, eliminating the need for an LLM to guess what is relevant.
      \end{exampleblock}
  \end{columns}
\end{frame}

% Slide 5
\begin{frame}{The pipeline bridges raw source code to a ranked context window}
  \begin{itemize}
    \item \textbf{Indexing Phase:} Parse the repository and build the dependency graph.
    \item \textbf{Retrieval Phase:} Select lexical seeds, run Personalized PageRank, and package context.
  \end{itemize}
  \vspace{10pt}
  \centering
  \includegraphics[width=0.9\textwidth,height=0.6\textheight,keepaspectratio]{/Users/alexandrutoma/dev/codegraph/thesis/assets/pipeline_architecture.png}
\end{frame}

% Slide 6
\begin{frame}{Behind the scenes: The Indexing Phase}
  \begin{block}{Offline Graph Building}
    The repository is parsed into a structural graph representation once per version, ensuring fast access during retrieval.
  \end{block}
  \vspace{10pt}
  \begin{itemize}
    \item \textbf{Entity Extraction:} Uses \HL{Tree-sitter} to parse a four-entity model: files, classes, functions, and methods.
    \item \textbf{Structural Connections:} Two-pass graph construction connects explicit \texttt{CONTAINS}, \texttt{IMPORTS}, \texttt{CALLS}, and \texttt{INHERITS\_FROM} edges.
    \item \textbf{Conservative Resolution:} Ambiguous call edges are dropped in favor of precision over false connectivity.
  \end{itemize}
\end{frame}

% Slide 7
\begin{frame}{Behind the scenes: The Retrieval Phase}
  \begin{exampleblock}{Online Execution (Ranking)}
    Each query dynamically activates seeds and propagates relevance through the graph deterministically.
  \end{exampleblock}
  \vspace{10pt}
  \begin{itemize}
    \item \textbf{Multi-Signal Seeding:} Combines Entity Matches, BM25 Keywords, and Issue Path Hints.
    \item \textbf{Deterministic Ranking:} Personalized PageRank relies on exact matrix multiplication. The \pos{same seeds produce the exact same ranking}.
    \item \textbf{Context Delivery:} Final ranked nodes undergo deduplication and are delivered within a strict 6,000-token budget.
  \end{itemize}
\end{frame}

% Slide 8
\begin{frame}{Global structural ranking drives the empirical gain}
  \begin{columns}[c]
    \column{0.45\textwidth}
      \textbf{Strong Empirical Gains:}
      \begin{itemize}
        \item CodeGraph reaches \pos{74.0\% Recall@10} on SWE-bench Lite.
        \item Improves upon the strongest lexical BM25 baseline by \HL{16.0 percentage points}.
      \end{itemize}
      \vspace{10pt}
      \textbf{Global vs Local:}
      \begin{itemize}
        \item A naive one-hop graph expansion only reaches 51.0\%, proving that gains come from \HL{global structural ranking}.
      \end{itemize}
      
    \column{0.55\textwidth}
      \centering
      \includegraphics[width=\textwidth,height=0.7\textheight,keepaspectratio]{/Users/alexandrutoma/dev/codegraph/thesis/assets/recall_trajectory.png}
  \end{columns}
\end{frame}

% Slide 9
\begin{frame}{The ultimate retrieval ceiling is seed reachability}
  \begin{columns}[c]
    \column{0.48\textwidth}
      \begin{alertblock}{Zero-Recall Failures (78 cases)}
        Retrieval quality is bounded by whether initial lexical signals can place probability mass into a structurally connected region.
      \end{alertblock}
      \vspace{10pt}
      \begin{itemize}
        \item \textbf{Connectivity Gaps (40 cases):} Disconnected graphs, dynamic imports.
        \item \textbf{Semantic Gaps (38 cases):} No vocabulary overlap.
      \end{itemize}
      
    \column{0.5\textwidth}
      \centering
      \includegraphics[width=\textwidth,height=0.65\textheight,keepaspectratio]{/Users/alexandrutoma/dev/codegraph/thesis/assets/failure_taxonomy.png}
  \end{columns}
\end{frame}

% Slide 10
\begin{frame}{Conclusion}
  \begin{block}{Scientific Takeaways}
    \begin{itemize}
      \item \textbf{Deterministic lexical-to-graph retrieval} delivers strong context effectively (74.0\% Recall@10).
      \item It maintains \pos{sub-second median latency}.
      \item The retrieval ceiling is determined by \pos{seed reachability}, not just ranking sophistication.
    \end{itemize}
  \end{block}
  \vspace{10pt}
  \begin{alertblock}{Important Boundaries}
    Static-analysis dependence, Python-only scope, and file-level evaluation remain core limitations.
  \end{alertblock}
\end{frame}

% Slide 11
\begin{frame}
  \centering
  \Huge \textbf{Thank you!} \\
  \vspace{20pt}
  
\begin{tikzpicture}[scale=1]
    \draw[thick, draw=negative, fill=yellow!80!orange] (0,0) circle (2);
    \fill[negative] ( -0.7, 0.5) circle (0.3);
    \fill[negative] (  0.7, 0.5) circle (0.3);
    \draw[thick, draw=negative, line width=3pt] (-1,-0.5) arc (180:360:1 and 0.8);
  \end{tikzpicture} \\
  \vspace{20pt}
  \Huge Questions?
\end{frame}

\end{document}
